\section{The review as a case study}
\label{sec:reflexive}

This review was performed by the class of system it studies. Rather than present the result
as though a human had produced it, we logged the process and analyse it here. Three logs were
kept contemporaneously and are released in full: a timeline, an error log, and a record of
every human intervention. Nothing was reconstructed after the fact and nothing was deleted.

\subsection{Two frontier models cannot agree on what an AI scientist is}

The most consequential finding concerns not the AI's competence but the protocol it wrote.
Layer-1 screening ran every candidate past two frontier models with different prompt framings.
They agreed on \rawAgree{} of \nScreened{} decisions ($\kappa = \kappaThree{}$), which by
convention is ``substantial'' agreement --- and which we initially read as a reassuring number.

Reading the disagreements instead of the summary statistic told a different story. The
\nAdjudicate{} contested decisions were not scattered; they concentrated on seven classes:

\begin{enumerate}\itemsep1pt
\item interactive ideation and writing aids, where the human makes the research judgments;
\item automated literature-review pipelines, which perform a research stage without an agent loop;
\item Bayesian-optimisation and robotic infrastructure for self-driving laboratories;
\item studies measuring human use of AI in writing, and detectors for AI-written text;
\item data-science agents whose tasks may or may not be scientific;
\item general-purpose ``deep research'' web agents;
\item institutional ethics frameworks for researchers using AI.
\end{enumerate}

Every one of these asks the same question --- \emph{what counts as an AI doing science?} ---
and the protocol, written by the AI and approved by the human, did not answer it. Two
independent frontier models disagreeing systematically rather than randomly is evidence about
the field's conceptual state, not merely about rater noise: this literature has not settled
its own boundaries, and a reviewer who screens it with a single pass will silently impose one
resolution among several defensible ones.

The remedy cost a full re-screen of the corpus (amendment A1). We note what did \emph{not}
surface the problem: the protocol document itself, human approval of that document, and the
headline agreement statistic. What surfaced it was redundancy at the definitional layer plus
the decision to inspect disagreements individually.

\subsection{An error taxonomy with measured rates}

Five failures were logged. We report them with rates because rates are what the literature
we mapped almost never provides.

\paragraph{Silent-empty success.} The first screening workflow returned success in 106
milliseconds having spawned zero agents: an argument arrived as a JSON string where the
orchestration code expected an object, so the job-construction loop produced an empty list.
The failure signal was a success signal. Detection required comparing the reported output
against an expected count.

\paragraph{Silent item-dropping in structured extraction.} Agents instructed to emit exactly
one record per input item occasionally emitted fewer, at a measured rate of 0.5\% in screening
(7 of 1{,}359) and 0.4\% in coding (3 of 811). Nothing in the output indicated the omission.
Undetected, this would have appeared as a slightly smaller corpus and nothing else. Detection
required per-item reconciliation against the input roster.

\paragraph{Reasoning over nothing.} A metadata-lookup defect passed two papers to a screening
step with empty titles and abstracts. The step duly produced confident-looking decisions about
them. The affected records were precisely the two that human expertise had rescued from the
search queries --- the pipeline silently nullified an earlier human correction. Detection came
from reading one flagged \texttt{UNDECIDED} justification rather than treating it as noise.

\paragraph{Query design defeated by an indexing detail.} An initial query returned 2{,}336
mostly irrelevant records because arXiv stems \texttt{auditable} to \texttt{audit*}. Notably,
the AI's first diagnosis of this anomaly --- boolean mis-parsing --- was wrong, and was
corrected by its own follow-up experiment. We record the wrong hypothesis as well as the right
one.

\paragraph{Protocol underdetermination.} Described above: 21.5\% of decisions left
undetermined by criteria that read as rigorous.

\subsection{What the human did}

The human co-author's interventions are logged individually. Three shaped the work materially.
He selected the research direction from four options the AI proposed; he set the publication
strategy, choosing dual-venue disclosure over the AI's recommendation; and at the scope
checkpoint he overruled the AI's boundary rule on deep-research agents, moving 183 papers into
the corpus. The AI had recommended and defended their exclusion. We think the human was right,
and record the disagreement rather than harmonising it.

He also declined the human screening audit that the protocol specified and the AI recommended.
That decision is logged too, and its cost is stated in Section~\ref{sec:limits}: it is the
single largest methodological weakness of this review.

\subsection{What this implies}

Three observations generalise beyond this project.

First, \textbf{every failure we logged was caught by a mechanical check that cost almost
nothing}: expected-versus-actual reconciliation, refusing to process empty input, and running
two raters against each other. None required sophisticated oversight. None would have been
visible in the finished paper.

Second, \textbf{summary statistics concealed the problems that individual records revealed}. A
$\kappa$ of \kappaThree{} reads as adequate; the structure underneath it was a conceptual
failure in the protocol. A workflow's success flag read as success. An item count read as
complete. In each case the aggregate was reassuring and the detail was not.

Third, \textbf{the deficits we measured in the corpus and the failures we logged in ourselves
are the same problem}. \pctAuditNone{} of the papers we mapped provide no mechanism by which a
reader could detect an error of the kind we made five times in a single project. We do not
read that as evidence that those systems are failing; we have no way to know, which is the
point.

\subsection{On the honesty of an AI reporting its own failures}

A caveat we cannot resolve from the inside. This section was written by the system whose
failures it reports. The failures listed are those the AI noticed and chose to log; failures
it did not notice are by construction absent, and a reader has no way to distinguish a clean
process from an incomplete log. We have tried to make the claim checkable rather than asking
for trust: the logs are timestamped in git history, the data files permit independent
recomputation of every statistic, and the failure modes named here are mechanically detectable
in the released artifacts. That is weaker than external audit, and we prefer to say so.
